\documentclass[12pt]{article}

% -------------------------------------------------------
% PACKAGES
% -------------------------------------------------------
\usepackage[margin=1in]{geometry}
\usepackage{amsmath,amssymb}
\usepackage{booktabs}
\usepackage{array}
\usepackage{multirow}
\usepackage{graphicx}
\usepackage{xcolor}
\usepackage[colorlinks=true,linkcolor=blue!60!black,urlcolor=blue!60!black,citecolor=blue!60!black,
  pdftitle={A Study Comparing NHL Player Salaries Across Positions and Geographic Regions Using Finite Population Sampling},
  pdfauthor={Neill Botushanski}]{hyperref}
\usepackage{setspace}
\usepackage{parskip}
\usepackage{titlesec}
\usepackage{fancyhdr}
\usepackage{enumitem}
\usepackage{float}
\usepackage{listings}
\usepackage{tcolorbox}
\usepackage{tabularx}
\usepackage{caption}
\usepackage{xurl}

% -------------------------------------------------------
% HEADER / FOOTER
% -------------------------------------------------------
\setlength{\headheight}{14.5pt}
\pagestyle{fancy}
\fancyhf{}
\lhead{\small MATH 3430 -- NHL Salary Analysis}
\rhead{\small York University, 2026}
\cfoot{\thepage}
\renewcommand{\headrulewidth}{0.4pt}

% -------------------------------------------------------
% SECTION FORMATTING
% -------------------------------------------------------
\titleformat{\section}{\large\bfseries}{\thesection}{1em}{}[\titlerule]
\titleformat{\subsection}{\normalsize\bfseries}{\thesubsection}{1em}{}
\titleformat{\subsubsection}{\normalsize\itshape}{\thesubsubsection}{1em}{}

% -------------------------------------------------------
% LISTINGS (R CODE)
% -------------------------------------------------------
\lstdefinestyle{Rcode}{
  language=R,
  basicstyle=\footnotesize\ttfamily,
  keywordstyle=\color{blue!70!black}\bfseries,
  commentstyle=\color{gray},
  stringstyle=\color{green!50!black},
  numbers=left,
  numberstyle=\tiny\color{gray},
  numbersep=8pt,
  breaklines=true,
  showstringspaces=false,
  frame=single,
  backgroundcolor=\color{gray!8},
  rulecolor=\color{gray!40},
  xleftmargin=2em,
  framexleftmargin=2em,
  aboveskip=6pt,
  belowskip=6pt,
  captionpos=b
}

% -------------------------------------------------------
% CUSTOM MATH COMMANDS
% -------------------------------------------------------
\newcommand{\VarHat}{\widehat{\mathrm{Var}}}
\newcommand{\SEHat}{\widehat{\mathrm{SE}}}

% -------------------------------------------------------
% DOCUMENT
% -------------------------------------------------------
\begin{document}

% ==============================================================
% TITLE PAGE
% ==============================================================
\begin{titlepage}
  \centering
  \vspace*{2cm}
  {\LARGE\bfseries A Study Comparing NHL Player Salaries\\[0.5em]
  Across Positions and Geographic Regions\\[0.5em]
  Using Finite Population Sampling\par}
  \vspace{1.5cm}
  {\large MATH 3430 -- Sample Survey Design\par}
  {\large York University\par}
  \vspace{1cm}
  \rule{0.5\textwidth}{0.4pt}\\[0.6cm]
  {\large Neill Botushanski\par}
  \vspace{0.3cm}
  {\normalsize Group project, completed with two teammates\par}
  \vspace{0.8cm}
  \rule{0.5\textwidth}{0.4pt}\\[0.6cm]
  {\large April 6, 2026\par}
  \vfill
  \begin{tcolorbox}[colback=gray!8, colframe=gray!50, width=0.85\textwidth]
    \small\centering
    This report presents a complete finite-population survey analysis of NHL player
    salaries using SRSWOR, Stratified Sampling with Proportional Allocation, and
    Stratified Sampling with Neyman Allocation. Strata are defined by playing position.
    All computations were performed in \textsf{R} with \texttt{set.seed(123)}.
  \end{tcolorbox}
\end{titlepage}

% ==============================================================
% TABLE OF CONTENTS
% ==============================================================
\tableofcontents
\newpage

% ==============================================================
\section{Abstract}
% ==============================================================

This study applies finite population sampling methods to estimate salary
characteristics of National Hockey League (NHL) players.
The population is $N = 787$ players, each with a recorded pre-tax annual salary and a
tax-adjusted after-tax salary, along with playing position (Forward, Defenseman,
Goaltender) and team country (Canada or United States).
Three sampling designs are evaluated at a fixed total sample size $n = 120$:
(1) Simple Random Sampling Without Replacement (SRSWOR),
(2) Stratified Sampling with Proportional Allocation, and
(3) Stratified Sampling with Neyman Allocation, where strata correspond to playing
position.
For each design and salary variable, we report the estimated population mean, its
variance, standard error, and 95\% confidence interval, together with the estimated
population total.
Relative efficiency comparisons confirm that both stratified designs outperform SRSWOR,
with Neyman allocation achieving the lowest variance.
A key substantive finding is that tax adjustment substantially reverses the
Canada-vs-United States salary ranking: U.S.\ teams have higher pre-tax average
salaries for defensemen and goaltenders, but after taxes Canadian defensemen and
forwards retain more, and the gap for goaltenders all but disappears.

\noindent\textbf{Keywords:} SRSWOR, stratified sampling, Neyman allocation,
proportional allocation, finite population correction, NHL salaries, relative efficiency.

\newpage

% ==============================================================
\section{Introduction}
\label{sec:intro}
% ==============================================================

Salary structures in professional sports leagues are shaped by player performance,
market size, contract length, and the tax environment of each team's home jurisdiction.
In the NHL, forwards, defensemen, and goaltenders fulfil fundamentally different
on-ice roles; whether these positional differences translate into systematic salary gaps
is an empirical question.
Furthermore, because tax burdens differ between Canadian and U.S.\ jurisdictions, the
pre-tax headline salary may misstate the real economic value of a contract.
Without adjusting for taxes, a naive comparison of nominal salaries could lead to
incorrect conclusions about which players are better compensated.

This project applies finite population sampling theory to address these questions
rigorously.
The 787-player dataset is treated as a finite population.
Rather than analysing the full population directly, we draw probability samples and
use them to construct estimators, thereby demonstrating how estimator precision depends
on the chosen sampling design.
This is precisely the setting of finite population inference: the population is fixed
and finite, randomness enters only through the sampling mechanism, and the goal is
to compare designs on the basis of the variance they induce.

\subsection*{Research Objectives}

\begin{enumerate}[nosep]
  \item Estimate the population mean NHL pre-tax salary $\mu_{\mathrm{pre}}$ and
        population mean after-tax salary $\mu_{\mathrm{after}}$, together with the
        corresponding population totals.
  \item Quantify estimator precision under three designs: SRSWOR, proportional
        stratified, and Neyman stratified sampling.
  \item Compare designs via relative efficiency (RE).
  \item Examine whether mean salary differs across the three playing positions.
  \item Assess how tax adjustment alters the Canada vs.\ United States salary
        comparison.
\end{enumerate}

\newpage

% ==============================================================
\section{Population and Data Description}
\label{sec:data}
% ==============================================================

\subsection{Population Frame}

The sampling frame is the complete set of $N = 787$ NHL players retained after the
following cleaning steps:
\begin{itemize}[nosep]
  \item Any row with a missing value in Team, Country, Player, Position, SalaryPreTax,
        or SalaryAfterTax was removed.
  \item The position code \texttt{C} (centre) was recoded to \texttt{F} (forward),
        yielding three mutually exclusive and exhaustive position categories.
  \item Rows with Position outside $\{F, D, G\}$ or Country outside $\{C, U\}$ were
        removed.
\end{itemize}

A derived variable $T_i = Y_i^{\mathrm{pre}} - Y_i^{\mathrm{after}}$ captures the
dollar amount lost to taxes for player $i$.

\subsection{Population Composition}

\begin{table}[H]
\centering
\caption{Population counts by playing position ($N = 787$)}
\label{tab:pop}
\begin{tabular}{lrr}
\toprule
Position & Count ($N_h$) & Proportion ($W_h = N_h/N$) \\
\midrule
Defenseman (D) & 269 & 0.3418 \\
Forward (F)    & 448 & 0.5693 \\
Goaltender (G) &  70 & 0.0889 \\
\midrule
Total          & 787 & 1.0000 \\
\bottomrule
\end{tabular}
\end{table}

\begin{table}[H]
\centering
\caption{Population counts by country}
\label{tab:pop_country}
\begin{tabular}{lrr}
\toprule
Country           & Count & Proportion \\
\midrule
Canada (C)        & 164   & 0.2084 \\
United States (U) & 623   & 0.7916 \\
\midrule
Total             & 787   & 1.0000 \\
\bottomrule
\end{tabular}
\end{table}

\subsection{Population-Level Salary Summaries}

Because the full dataset is available, the true population means serve as benchmarks
against which the sample estimates can be evaluated.

\begin{table}[H]
\centering
\caption{Population mean pre-tax salary by position and region}
\label{tab:pop_pretax}
\begin{tabular}{llrrr}
\toprule
Position & Region & Count & Mean Pre-Tax (\$) & SD (\$) \\
\midrule
D & Canada        &  52 & 3{,}545{,}124 & 2{,}757{,}696 \\
D & United States & 217 & 3{,}698{,}438 & 2{,}899{,}817 \\
F & Canada        &  97 & 3{,}805{,}874 & 2{,}887{,}633 \\
F & United States & 351 & 3{,}585{,}417 & 2{,}796{,}670 \\
G & Canada        &  15 & 3{,}195{,}444 & 2{,}658{,}782 \\
G & United States &  55 & 3{,}826{,}136 & 2{,}843{,}115 \\
\midrule
Overall           & & 787 & 3{,}650{,}480 & --- \\
\bottomrule
\end{tabular}
\end{table}

\begin{table}[H]
\centering
\caption{Population mean after-tax salary by position and region}
\label{tab:pop_aftertax}
\begin{tabular}{llrrr}
\toprule
Position & Region & Count & Mean After-Tax (\$) & SD (\$) \\
\midrule
D & Canada        &  52 & 1{,}823{,}065 & 1{,}415{,}306 \\
D & United States & 217 & 1{,}595{,}217 & 1{,}287{,}632 \\
F & Canada        &  97 & 1{,}963{,}960 & 1{,}510{,}241 \\
F & United States & 351 & 1{,}532{,}326 & 1{,}196{,}062 \\
G & Canada        &  15 & 1{,}657{,}259 & 1{,}381{,}689 \\
G & United States &  55 & 1{,}677{,}503 & 1{,}297{,}076 \\
\midrule
Overall           & & 787 & 1{,}634{,}604 & --- \\
\bottomrule
\end{tabular}
\end{table}

\begin{table}[H]
\centering
\caption{Population mean tax lost (\$) by position and region}
\label{tab:pop_taxlost}
\begin{tabular}{llrr}
\toprule
Position & Region & Count & Mean Tax Lost (\$) \\
\midrule
D & Canada        &  52 & 1{,}722{,}059 \\
D & United States & 217 & 2{,}103{,}220 \\
F & Canada        &  97 & 1{,}841{,}914 \\
F & United States & 351 & 2{,}053{,}091 \\
G & Canada        &  15 & 1{,}538{,}186 \\
G & United States &  55 & 2{,}148{,}633 \\
\midrule
Overall           & & 787 & 2{,}015{,}876 \\
\bottomrule
\end{tabular}
\end{table}

\noindent The overall population mean pre-tax salary is \$3{,}650{,}480 and the mean
after-tax salary is \$1{,}634{,}604, implying that the average player loses roughly
\$2{,}015{,}876 -- approximately 55\% of gross salary -- to taxes.
Crucially, the regional ranking flips after taxes: for forwards and defensemen,
Canadian players retain \emph{more} after-tax income despite having pre-tax salaries
that are comparable to, or in some cases lower than, those of their U.S.\ counterparts.

\newpage

% ==============================================================
\section{Sampling Design and Methods}
\label{sec:methods}
% ==============================================================

\subsection{Notation}

Let $U = \{1, 2, \ldots, N\}$ be the population of $N = 787$ units and $s$ a
probability sample of size $n = 120$.
The overall sampling fraction is
\[
  f = \frac{n}{N} = \frac{120}{787} \approx 0.1525.
\]
For stratified designs we index strata by $h \in \{D, F, G\}$, with $N_h$ the stratum
population size, $n_h$ the stratum sample size, $W_h = N_h / N$ the stratum weight,
and $f_h = n_h / N_h$ the stratum sampling fraction.
Let $\bar{Y}_h$ denote the true stratum mean and $S_h^2$ the true stratum variance,
with sample analogues $\bar{y}_h$ and $s_h^2$.

\subsection{Simple Random Sampling Without Replacement (SRSWOR)}
\label{sec:srs}

Under SRSWOR every subset of size $n$ from $U$ has equal probability of selection.
The estimator of the population mean is the sample mean:
\[
  \hat{\bar{Y}}_{\mathrm{SRS}} = \bar{y} = \frac{1}{n}\sum_{i \in s} y_i.
\]
Its estimated variance incorporates the finite population correction (FPC):
\[
  \VarHat\!\left(\hat{\bar{Y}}_{\mathrm{SRS}}\right)
  = \underbrace{\left(1 - \frac{n}{N}\right)}_{\text{FPC}} \frac{s^2}{n},
  \quad \text{where } s^2 = \frac{1}{n-1}\sum_{i \in s}(y_i - \bar{y})^2.
\]
The FPC reduces the variance below the with-replacement formula by a factor
$(1-f)$; here $1 - f \approx 0.8475$.

The estimated population total and its standard error are:
\[
  \hat{T}_{\mathrm{SRS}} = N\,\hat{\bar{Y}}_{\mathrm{SRS}},
  \qquad
  \SEHat\!\left(\hat{T}_{\mathrm{SRS}}\right) = N\,\SEHat\!\left(\hat{\bar{Y}}_{\mathrm{SRS}}\right).
\]

A 95\% confidence interval for $\bar{Y}$ is constructed as:
\[
  \hat{\bar{Y}}_{\mathrm{SRS}}
  \;\pm\; 1.96\;\SEHat\!\left(\hat{\bar{Y}}_{\mathrm{SRS}}\right).
\]

\subsection{Stratified Sampling with Proportional Allocation}
\label{sec:prop}

The population is partitioned into $H = 3$ strata by playing position.
Under proportional allocation the stratum sample sizes are
\[
  n_h^{\mathrm{prop}} = n \cdot W_h = n \cdot \frac{N_h}{N}, \qquad h \in \{D,F,G\}.
\]
Since $n_h^{\mathrm{prop}}$ is generally non-integer, the \emph{largest-remainder
method} is applied: assign $\lfloor n_h^{\mathrm{prop}} \rfloor$ to each stratum,
then distribute the remaining $n - \sum_h \lfloor n_h^{\mathrm{prop}} \rfloor$
units one by one to the strata with the largest fractional remainders, subject to
$n_h \leq N_h$. The resulting integer allocation is shown in Table~\ref{tab:prop_alloc}.

\begin{table}[H]
\centering
\caption{Proportional allocation by position}
\label{tab:prop_alloc}
\begin{tabular}{lrrrr}
\toprule
Position & $N_h$ & $W_h$ & $n_h^{\text{raw}} = nW_h$ & $n_h$ (integer) \\
\midrule
D & 269 & 0.3418 & 41.02 & 41 \\
F & 448 & 0.5693 & 68.32 & 68 \\
G &  70 & 0.0889 & 10.67 & 11 \\
\midrule
Total & 787 & 1.0000 & 120.00 & 120 \\
\bottomrule
\end{tabular}
\end{table}

Within each stratum an independent SRSWOR of size $n_h$ is drawn.
The stratified estimator of the population mean is:
\[
  \hat{\bar{Y}}_{\mathrm{st}} = \sum_{h=1}^{H} W_h\,\bar{y}_h.
\]
Its estimated variance is:
\[
  \VarHat\!\left(\hat{\bar{Y}}_{\mathrm{st}}\right)
  = \sum_{h=1}^{H} W_h^2 \left(1 - \frac{n_h}{N_h}\right)\frac{s_h^2}{n_h}.
\]
Each term contains a stratum-level FPC $(1 - n_h/N_h)$ and is weighted by $W_h^2$.
The estimated population total is $\hat{T}_{\mathrm{st}} = N\,\hat{\bar{Y}}_{\mathrm{st}}$.

\subsection{Stratified Sampling with Neyman Allocation}
\label{sec:neyman}

Neyman (optimal) allocation minimises $\mathrm{Var}(\hat{\bar{Y}}_{\mathrm{st}})$
for a fixed total sample size $n$:
\[
  n_h^{\mathrm{Ney}} = n \cdot \frac{N_h S_h}{\displaystyle\sum_{k=1}^{H} N_k S_k}.
\]
Strata with larger population size $N_h$ or greater variability $S_h$ receive more
sample units. Because the true $S_h$ are unknown in practice, we substitute the
population standard deviations computed from the full dataset (oracle allocation).
Since $S_h$ differs between pre-tax and after-tax salary, separate Neyman allocations
are computed for each variable.

\subsubsection*{Neyman Allocation for Pre-Tax Salary}

Population stratum standard deviations:
$S_D = 2{,}868{,}489$, $S_F = 2{,}814{,}813$, $S_G = 2{,}797{,}913$.

\begin{align*}
  \sum_h N_h S_h
  &= (269)(2{,}868{,}489) + (448)(2{,}814{,}813) + (70)(2{,}797{,}913) \\
  &= 771{,}623{,}582 + 1{,}261{,}036{,}342 + 195{,}853{,}919 \\
  &= 2{,}228{,}513{,}843.
\end{align*}
(The products are computed from the unrounded $S_h$, so they differ slightly from
the product of the rounded values shown.)
\begin{align*}
  n_D^{\text{raw}} &= 120 \times \frac{771{,}623{,}582}{2{,}228{,}513{,}843} = 41.550, \\[4pt]
  n_F^{\text{raw}} &= 120 \times \frac{1{,}261{,}036{,}342}{2{,}228{,}513{,}843} = 67.904, \\[4pt]
  n_G^{\text{raw}} &= 120 \times \frac{195{,}853{,}919}{2{,}228{,}513{,}843} = 10.546.
\end{align*}
The integer parts (41, 67, 10) sum to 118, so the largest-remainder method assigns the
two remaining units to F and D, giving $n_D = 42$, $n_F = 68$, $n_G = 10$.

\begin{table}[H]
\centering
\caption{Neyman allocation for pre-tax salary}
\label{tab:neyman_pre}
\begin{tabular}{lrrrrrr}
\toprule
Position & $N_h$ & $S_h$ & $W_h$ & $N_h S_h$ & $n_h^{\text{raw}}$ & $n_h$ \\
\midrule
D & 269 & 2{,}868{,}489 & 0.342 & 771{,}623{,}582   & 41.550 & 42 \\
F & 448 & 2{,}814{,}813 & 0.569 & 1{,}261{,}036{,}342 & 67.904 & 68 \\
G &  70 & 2{,}797{,}913 & 0.089 & 195{,}853{,}919   & 10.546 & 10 \\
\midrule
Total & 787 & & & 2{,}228{,}513{,}843 & 120.000 & 120 \\
\bottomrule
\end{tabular}
\end{table}

\subsubsection*{Neyman Allocation for After-Tax Salary}

Population stratum standard deviations:
$S_D = 1{,}313{,}623$, $S_F = 1{,}281{,}268$, $S_G = 1{,}305{,}405$.

\begin{align*}
  \sum_h N_h S_h
  &= (269)(1{,}313{,}623) + (448)(1{,}281{,}268) + (70)(1{,}305{,}405) \\
  &= 353{,}364{,}587 + 574{,}008{,}064 + 91{,}378{,}350 \\
  &= 1{,}018{,}751{,}001.
\end{align*}
\begin{align*}
  n_D^{\text{raw}} &= 120 \times \frac{353{,}364{,}587}{1{,}018{,}751{,}001} = 41.623, \\[4pt]
  n_F^{\text{raw}} &= 120 \times \frac{574{,}008{,}064}{1{,}018{,}751{,}001} = 67.613, \\[4pt]
  n_G^{\text{raw}} &= 120 \times \frac{91{,}378{,}350}{1{,}018{,}751{,}001} = 10.764.
\end{align*}
The integer parts (41, 67, 10) again sum to 118; the two remaining units go to G and D,
giving $n_D = 42$, $n_F = 67$, $n_G = 11$.

\begin{table}[H]
\centering
\caption{Neyman allocation for after-tax salary}
\label{tab:neyman_after}
\begin{tabular}{lrrrrrr}
\toprule
Position & $N_h$ & $S_h$ & $W_h$ & $N_h S_h$ & $n_h^{\text{raw}}$ & $n_h$ \\
\midrule
D & 269 & 1{,}313{,}623 & 0.342 & 353{,}364{,}587   & 41.623 & 42 \\
F & 448 & 1{,}281{,}268 & 0.569 & 574{,}008{,}064   & 67.613 & 67 \\
G &  70 & 1{,}305{,}405 & 0.089 &  91{,}378{,}350   & 10.764 & 11 \\
\midrule
Total & 787 & & & 1{,}018{,}751{,}001 & 120.000 & 120 \\
\bottomrule
\end{tabular}
\end{table}

\subsection{Relative Efficiency}
\label{sec:re}

The relative efficiency of design $A$ relative to design $B$ is defined as:
\[
  \mathrm{RE}(A \text{ vs.\ } B)
  = \frac{\VarHat\!\left(\hat{\bar{Y}}_{B}\right)}{\VarHat\!\left(\hat{\bar{Y}}_{A}\right)}.
\]
A value $\mathrm{RE} > 1$ means design $A$ has smaller estimated variance (is more
efficient) than design $B$.

\newpage

% ==============================================================
\section{Results}
\label{sec:results}
% ==============================================================

\subsection{Pre-Tax Salary: SRSWOR}

A SRSWOR of $n = 120$ was drawn with \texttt{set.seed(123)}.

\[
  \hat{\bar{Y}}_{\mathrm{SRS}}^{\mathrm{pre}} = \$4{,}170{,}330,
  \qquad
  s^2 = 8{,}925{,}643{,}801{,}435.
\]

Estimated variance with FPC:
\[
  \VarHat\!\left(\hat{\bar{Y}}_{\mathrm{SRS}}^{\mathrm{pre}}\right)
  = \left(1 - \frac{120}{787}\right)\frac{8{,}925{,}643{,}801{,}435}{120}
  = (0.8475)(74{,}380{,}365{,}012)
  = 63{,}039{,}013{,}295.
\]

\[
  \SEHat = \sqrt{63{,}039{,}013{,}295} = \$251{,}076,
  \qquad
  \text{95\% CI: } (\$3{,}678{,}222,\ \$4{,}662{,}439).
\]

\[
  \hat{T}_{\mathrm{SRS}}^{\mathrm{pre}} = 787 \times 4{,}170{,}330 = \$3{,}282{,}049{,}985,
  \qquad
  \SEHat(\hat{T}) = 787 \times 251{,}076 = \$197{,}596{,}812.
\]

\subsection{Pre-Tax Salary: Stratified Proportional}

\[
  \hat{\bar{Y}}_{\mathrm{st,prop}}^{\mathrm{pre}} = \$3{,}362{,}231,
  \qquad
  \VarHat = 54{,}651{,}670{,}893,
  \qquad
  \SEHat = \$233{,}777.
\]
\[
  \text{95\% CI: } (\$2{,}904{,}028,\ \$3{,}820{,}434).
\]
\[
  \hat{T}_{\mathrm{st,prop}}^{\mathrm{pre}} = \$2{,}646{,}076{,}038,
  \qquad
  \SEHat(\hat{T}) = \$183{,}982{,}474.
\]

\subsection{Pre-Tax Salary: Stratified Neyman}

\[
  \hat{\bar{Y}}_{\mathrm{st,Ney}}^{\mathrm{pre}} = \$3{,}349{,}263,
  \qquad
  \VarHat = 53{,}488{,}449{,}024,
  \qquad
  \SEHat = \$231{,}276.
\]
\[
  \text{95\% CI: } (\$2{,}895{,}962,\ \$3{,}802{,}563).
\]
\[
  \hat{T}_{\mathrm{st,Ney}}^{\mathrm{pre}} = \$2{,}635{,}869{,}698,
  \qquad
  \SEHat(\hat{T}) = \$182{,}013{,}975.
\]

\subsection{Pre-Tax Salary Summary}

\begin{table}[H]
\centering
\caption{Pre-tax salary estimator comparison}
\label{tab:pre_compare}
\begin{tabular}{lrrrrr}
\toprule
Method & $\hat{\bar{Y}}$ (\$) & SE (\$) & CI Lower (\$) & CI Upper (\$) & $\hat{T}$ (\$) \\
\midrule
SRSWOR       & 4{,}170{,}330 & 251{,}076 & 3{,}678{,}222 & 4{,}662{,}439 & 3{,}282{,}049{,}985 \\
Proportional & 3{,}362{,}231 & 233{,}777 & 2{,}904{,}028 & 3{,}820{,}434 & 2{,}646{,}076{,}038 \\
Neyman       & 3{,}349{,}263 & 231{,}276 & 2{,}895{,}962 & 3{,}802{,}563 & 2{,}635{,}869{,}698 \\
\bottomrule
\end{tabular}
\end{table}

Relative efficiencies for pre-tax salary:
\[
  \mathrm{RE}(\text{Prop vs SRS}) = \frac{63{,}039{,}013{,}295}{54{,}651{,}670{,}893} = 1.1535,
\]
\[
  \mathrm{RE}(\text{Neyman vs SRS}) = \frac{63{,}039{,}013{,}295}{53{,}488{,}449{,}024} = 1.1786,
\]
\[
  \mathrm{RE}(\text{Neyman vs Prop}) = \frac{54{,}651{,}670{,}893}{53{,}488{,}449{,}024} = 1.0217.
\]

The variance reduction implied by a relative efficiency is $1 - 1/\mathrm{RE}$.
Proportional allocation therefore achieves a 13.3\% variance reduction over SRSWOR and
Neyman allocation a 15.2\% reduction, an incremental 2.1\% beyond proportional.

\subsection{After-Tax Salary: SRSWOR}

The same SRSWOR sample is used for after-tax salary:
\[
  \hat{\bar{Y}}_{\mathrm{SRS}}^{\mathrm{after}} = \$1{,}899{,}015,
  \qquad
  s^2 = 1{,}993{,}664{,}126{,}684.
\]
\[
  \VarHat\!\left(\hat{\bar{Y}}_{\mathrm{SRS}}^{\mathrm{after}}\right)
  = (0.8475)\frac{1{,}993{,}664{,}126{,}684}{120}
  = 14{,}080{,}622{,}326,
  \qquad
  \SEHat = \$118{,}662.
\]
\[
  \text{95\% CI: } (\$1{,}666{,}438,\ \$2{,}131{,}592).
\]
\[
  \hat{T}_{\mathrm{SRS}}^{\mathrm{after}} = \$1{,}494{,}525{,}067,
  \qquad
  \SEHat(\hat{T}) = \$93{,}386{,}835.
\]

\subsection{After-Tax Salary: Stratified Proportional}

\[
  \hat{\bar{Y}}_{\mathrm{st,prop}}^{\mathrm{after}} = \$1{,}540{,}873,
  \qquad
  \VarHat = 12{,}518{,}558{,}006,
  \qquad
  \SEHat = \$111{,}886.
\]
\[
  \text{95\% CI: } (\$1{,}321{,}575,\ \$1{,}760{,}170).
\]
\[
  \hat{T}_{\mathrm{st,prop}}^{\mathrm{after}} = \$1{,}212{,}666{,}704,
  \qquad
  \SEHat(\hat{T}) = \$88{,}054{,}567.
\]

\subsection{After-Tax Salary: Stratified Neyman}

\[
  \hat{\bar{Y}}_{\mathrm{st,Ney}}^{\mathrm{after}} = \$1{,}535{,}275,
  \qquad
  \VarHat = 12{,}296{,}975{,}180,
  \qquad
  \SEHat = \$110{,}892.
\]
\[
  \text{95\% CI: } (\$1{,}317{,}927,\ \$1{,}752{,}623).
\]
\[
  \hat{T}_{\mathrm{st,Ney}}^{\mathrm{after}} = \$1{,}208{,}261{,}463,
  \qquad
  \SEHat(\hat{T}) = \$87{,}271{,}789.
\]

\subsection{After-Tax Salary Summary}

\begin{table}[H]
\centering
\caption{After-tax salary estimator comparison}
\label{tab:after_compare}
\begin{tabular}{lrrrrr}
\toprule
Method & $\hat{\bar{Y}}$ (\$) & SE (\$) & CI Lower (\$) & CI Upper (\$) & $\hat{T}$ (\$) \\
\midrule
SRSWOR       & 1{,}899{,}015 & 118{,}662 & 1{,}666{,}438 & 2{,}131{,}592 & 1{,}494{,}525{,}067 \\
Proportional & 1{,}540{,}873 & 111{,}886 & 1{,}321{,}575 & 1{,}760{,}170 & 1{,}212{,}666{,}704 \\
Neyman       & 1{,}535{,}275 & 110{,}892 & 1{,}317{,}927 & 1{,}752{,}623 & 1{,}208{,}261{,}463 \\
\bottomrule
\end{tabular}
\end{table}

Relative efficiencies for after-tax salary:
\[
  \mathrm{RE}(\text{Prop vs SRS}) = \frac{14{,}080{,}622{,}326}{12{,}518{,}558{,}006} = 1.1248,
\]
\[
  \mathrm{RE}(\text{Neyman vs SRS}) = \frac{14{,}080{,}622{,}326}{12{,}296{,}975{,}180} = 1.1450,
\]
\[
  \mathrm{RE}(\text{Neyman vs Prop}) = \frac{12{,}518{,}558{,}006}{12{,}296{,}975{,}180} = 1.0180.
\]

\subsection{Consolidated Efficiency Comparison}

\begin{table}[H]
\centering
\caption{Relative efficiency summary for both salary variables}
\label{tab:re_summary}
\begin{tabular}{lrr}
\toprule
Comparison & Pre-Tax RE & After-Tax RE \\
\midrule
Proportional vs.\ SRSWOR & 1.1535 & 1.1248 \\
Neyman vs.\ SRSWOR       & 1.1786 & 1.1450 \\
Neyman vs.\ Proportional & 1.0217 & 1.0180 \\
\bottomrule
\end{tabular}
\end{table}

\noindent Across both salary variables, stratified sampling consistently improves
precision.
The incremental gain from Neyman over proportional allocation is small (roughly 2\%
variance reduction) because the three position strata have nearly identical
within-stratum standard deviations.
When all $S_h$ are equal, Neyman allocation reduces exactly to proportional allocation,
and no further gains are possible from optimal within-stratum weighting.

\subsection{Interpretation of Substantive Results}

\subsubsection*{Salary distributions}

Both pre-tax and after-tax salary distributions are strongly right-skewed.
The majority of players earn near the league minimum (\$775{,}000), while a small
number of elite players command contracts up to \$17{,}000{,}000.
This concentration of extreme values drives the large within-stratum standard
deviations (around \$2{,}800{,}000 for pre-tax across all positions) and limits the
precision gains achievable through stratification by position alone.

\subsubsection*{Position comparisons}

Mean pre-tax salaries are remarkably similar across positions:
\$3{,}668{,}801 (D), \$3{,}633{,}150 (F), \$3{,}690{,}988 (G).
The same pattern holds for after-tax salaries.
Playing position is therefore a weak predictor of individual salary relative to
within-position variability, which explains the modest efficiency gains from
stratification.

\subsubsection*{Regional comparison and the tax reversal}

On a pre-tax basis, U.S.\ teams pay more than Canadian teams for defensemen
(\$3{,}698{,}438 vs.\ \$3{,}545{,}124) and goaltenders
(\$3{,}826{,}136 vs.\ \$3{,}195{,}444),
while Canadian teams pay more for forwards
(\$3{,}805{,}874 vs.\ \$3{,}585{,}417).
After taxes this picture changes substantially:
Canadian defensemen retain \$1{,}823{,}065 vs.\ \$1{,}595{,}217 for U.S.\ defensemen,
and Canadian forwards retain \$1{,}963{,}960 vs.\ \$1{,}532{,}326 for U.S.\ forwards.
Only for goaltenders does the difference effectively vanish
(\$1{,}657{,}259 Canada vs.\ \$1{,}677{,}503 United States).

The mean tax lost by U.S.\ players (\$2{,}103{,}220 D, \$2{,}053{,}091 F,
\$2{,}148{,}633 G) consistently exceeds that of Canadian players
(\$1{,}722{,}059 D, \$1{,}841{,}914 F, \$1{,}538{,}186 G),
confirming that higher U.S.\ nominal salaries are partly or fully offset by higher
tax burdens.

\newpage

% ==============================================================
\section{Conclusions and Discussion}
\label{sec:conclusions}
% ==============================================================

\subsection{Summary of Main Findings}

This study estimated NHL salary characteristics from a probability sample of $n = 120$
players drawn from a finite population of $N = 787$.

The true population mean pre-tax salary is \$3{,}650{,}480 and the true population
mean after-tax salary is \$1{,}634{,}604.
Both stratified estimators (proportional and Neyman) fell closer to these true values
than the SRSWOR estimator, which was inflated by high-salary outliers captured in the
particular SRSWOR draw.
This illustrates the importance of design: a more precise design is less susceptible
to chance imbalance in how outliers are represented.

Stratified sampling with positions as strata improves precision over SRSWOR in all
cases.
For pre-tax salary, proportional allocation reduces variance by 13.3\% and Neyman
allocation by 15.2\%. For after-tax salary, the figures are 11.1\% and 12.7\%.
The additional gain from Neyman over proportional is approximately 2\% in both cases,
reflecting near-uniform within-stratum variability across positions.

The most substantive finding is the regional tax reversal.
Pre-tax salary comparisons suggest U.S.\ teams offer higher compensation for
defensemen and goaltenders, yet after-tax salary comparisons show that Canadian
defensemen and forwards actually retain more income.
Any analysis of NHL compensation that relies solely on pre-tax salaries therefore
risks misleading conclusions about the financial incentives for players choosing
between Canadian and U.S.\ teams.

\subsection{Limitations}

\begin{enumerate}
  \item \textbf{Right-skewed distributions.} The high within-stratum variance driven
        by extreme salaries limits the efficiency gains from stratification by position.
        Stratifying on contract tier, age, or years of experience might produce more
        homogeneous strata and larger variance reductions.

  \item \textbf{Cross-sectional scope.} The dataset reflects a single season.
        Collective bargaining agreements and roster changes evolve over time; the
        findings may not generalise to other seasons.

  \item \textbf{Tax simplification.} The after-tax salary variable is based on
        simplified tax calculations and does not capture the full complexity of athlete
        taxation, including state or provincial rates, endorsement income, deferred
        compensation, or signing bonuses.

  \item \textbf{Oracle Neyman allocation.} Using full-population standard deviations
        to set Neyman allocation is an oracle procedure unavailable in a real survey
        setting, where $S_h$ must be estimated from a pilot sample or auxiliary data.
        This introduces additional uncertainty not quantified in the present analysis.

  \item \textbf{Single sample realisation.} All results are based on a single sample
        draw per design. A simulation study repeating the sampling many times would
        provide empirical confirmation of the relative efficiency figures.
\end{enumerate}

\newpage

% ==============================================================
\section*{References}
\addcontentsline{toc}{section}{References}
% ==============================================================

\begin{enumerate}
  \item Cochran, W.~G. (1977). \textit{Sampling Techniques} (3rd ed.). John Wiley \& Sons.
  \item Thompson, S.~K. (2012). \textit{Sampling} (3rd ed.). John Wiley \& Sons.
  \item Spotrac NHL Contract Average Dataset. Available at:
        \url{https://www.spotrac.com/nhl/rankings/player/_/year/2025/sort/contract_average}.
  \item The Commute Sports. NHL Salary Analysis (2019). Available at:
        \url{https://thecommutesports.com/2019/09/07/nhl-salary-analysis/}.
  \item R Core Team (2024). \textit{R: A Language and Environment for Statistical
        Computing}. R Foundation for Statistical Computing, Vienna, Austria.
        \url{https://www.R-project.org/}.
\end{enumerate}

\newpage

% ==============================================================
\appendix
\section*{Appendix A: Complete Annotated R Code}
\addcontentsline{toc}{section}{Appendix A: Complete Annotated R Code}
% ==============================================================

% The listing is read straight from the script so the two cannot drift apart.
\lstinputlisting[style=Rcode, caption={Complete R script for NHL salary analysis (MATH 3430)}]{../R/nhl_salary_sampling.R}

\end{document}
